% Draft 2. Section-specific decisions are recorded in ../SECTION-DECISIONS.md.
% Abstract-level method detail; short original evaluation; final-report framing.
% Adoption figures checked against GitHub API in October 2026; no citation count.
% House-price example stays here; the eventual single-paragraph abstract will summarise findings.
\section{Paper Description and Justification}
\label{sec:paper}

A house-price forecast of \$800{,}000 does not describe its uncertainty.
Two 95\% intervals, from \$780{,}000 to \$820{,}000 and from \$500{,}000 to \$1.1~million, support different decisions despite sharing that point estimate.
Probabilistic regression predicts a distribution over the outcome, conditional on the inputs, so uncertainty can vary between observations.
Its fit is evaluated on held-out data.

\citet{duan2020ngboost} introduce Natural Gradient Boosting, or NGBoost, for this purpose.\footnote{\url{https://proceedings.mlr.press/v119/duan20a.html}}
The method extends gradient boosting, which combines a sequence of weak learners, to estimate several parameters of a predictive distribution together.
With the Normal distribution used in this project, these parameters describe an input-dependent mean $\mu(x)$ and standard deviation $\sigma(x)$; the implementation represents the latter through $\log\sigma(x)$.
Fitting the parameters together can produce poorly balanced updates under an ordinary gradient.
For the logarithmic score, the natural gradient rescales the loss gradient using the inverse Fisher information, accounting for how parameter changes alter the distribution \cite{amari1998natural}.
The paper's contribution is a boosting framework that combines this update with a chosen distribution, base learner and proper scoring rule.

\textbf{Original evaluation.}
The paper compares NGBoost with probabilistic prediction methods including neural-network approaches, Gaussian processes and distributional regression, and separately tests its components through ablations \cite{duan2020ngboost}.
Its regression experiments use ten benchmark datasets, Normal predictions and depth-three decision trees.
The main measure is mean test negative log-likelihood (NLL), which evaluates the predictive distribution rather than only its mean; lower values are better.
For Yacht and Kin8nm, the published protocol repeats random 90/10 training/test holdouts twenty times.
Within each training portion, validation log-likelihood selects the number of boosting stages before refitting on the full training portion.

The original paper's ablation table provides the comparison investigated here.
On Yacht, the reported mean NLL is $0.20 \pm 0.26$ for NGBoost, against $3.29 \pm 0.20$ for ordinary-gradient multiparameter boosting and $2.02 \pm 0.21$ for the homoscedastic model, whose spread is constant across inputs.
The quoted spreads are standard errors over twenty repetitions.
On Kin8nm, ordinary-gradient boosting reports $-0.52 \pm 0.03$, NGBoost $-0.49 \pm 0.02$, and the homoscedastic model $-0.37 \pm 0.05$.
The ranges for the first two overlap, so the paper's own reporting criterion treats them as tied.
These comparisons are the claims this project tests, and they already show that NGBoost does not lead on every dataset.
This project examines those three variants on both datasets and extends the comparison to five further public datasets and constructed data with known conditional noise.
It does not reproduce the second-order ablation or the full comparison against external model families.

\textbf{Why this paper.}
The choice rests on publication quality, reproducibility and a directly testable contribution.
NGBoost appeared at the International Conference on Machine Learning in 2020 \cite{duan2020ngboost}.
ICML has an A\textsuperscript{*} rating in ICORE2026, as recorded in the project's August 2026 venue check \cite{icore2026icml}.
Venue standing is evidence about peer review, and reproducibility is a separate property that this project measures directly.

The official implementation is public under the Apache~2.0 licence \cite{ngboostrepo}.
This project runs that implementation at a recorded commit, with its own scripts for the repair and all experiments.\footnote{\label{fn:repo}Project repository: \url{https://github.com/Sakshamb13579/Replicating-Natural-Gradient-Boosting}}
Public data and code make the comparison inspectable.
A further criterion was that the released code ran to completion with only minor repairs: two missing import statements, which left the model, splits and stage selection unchanged (Section~\ref{sec:replication}).
Running the published implementation, rather than reimplementing it, avoids introducing changes that could alter its results.
The project uses scripted data acquisition and preparation with checksums, a synthetic data generator, repeated computational experiments, human annotation with an agreement measure, an LLM-as-judge evaluation and LaTeX reporting.
The NGBoost repository's 1{,}893 stars and 257 forks in October 2026 provide supporting evidence of community interest.

Finally, the ablations connect a specific algorithmic change to distributional performance.
Comparing natural and ordinary gradients tests the update rule, while the fixed-spread comparator examines the value of allowing uncertainty to vary with the input.
Constructed data makes the second comparison especially informative because the generating spread is known.
